\clearpage\corrigesChapitre{2}{Identifier la matière}
\corrigesSouples
\begin{blocLivre}{cor-2-1}\begin{corrige}{1}{Trier des échantillons}
\begin{questions}
\item Cuivre : corps pur ; eau sucrée : mélange homogène ; eau-huile : mélange hétérogène après décantation.
\item Non, le sucre dissous reste présent même si on ne voit pas ses entités.
\end{questions}
\end{corrige}\end{blocLivre}
\begin{blocLivre}{cor-2-2}\begin{corrige}{2}{Mesurer une masse volumique}
\schema[corrige]{densite}
\begin{questions}
\item $m=58,4-42,6=\qty{15,8}{g}$.
\item $\rho=15,8/20,0=\qty{0,790}{g.mL^{-1}}$.
\end{questions}
\end{corrige}\end{blocLivre}
\begin{blocLivre}{cor-2-3}\begin{corrige}{3}{Reconnaître un gaz}
\begin{questions}
\item Le dioxyde de carbone, \ce{CO2}.
\item Le dioxygène ravive une bûchette incandescente, dans un test encadré au laboratoire.
\end{questions}
\end{corrige}\end{blocLivre}
\begin{blocLivre}{cor-2-4}\begin{corrige}{4}{Air dans une pièce}
\begin{questions}
\item $V(\ce{O2})=0,21\times50=\qty{10,5}{m^3}$.
\item Le diazote représente environ 78 \% du volume.
\end{questions}
\end{corrige}\end{blocLivre}
\begin{blocLivre}{cor-2-5}\begin{corrige}{5}{Composition d’un alliage}
\begin{questions}
\item $100\times75,0/250=30,0\,\%$.
\item $m_{\mathrm{Cu}}=175\unit{g}$ et $w_{\mathrm{Cu}}=70,0\,\%$. Les pourcentages somment à 100.
\end{questions}
\end{corrige}\end{blocLivre}
\begin{blocLivre}{cor-2-6}\begin{corrige}{6}{Une piste de chromatographie}
\schema[corrige]{ccm-identification}
\begin{questions}
\item Les taches de M correspondent aux hauteurs de R et S : les deux espèces sont compatibles avec M.
\item Non, deux taches distinctes sont détectées dans la piste M.
\end{questions}
\end{corrige}\end{blocLivre}
\begin{blocLivre}{cor-2-7}\begin{corrige}{7}{Choisir un matériau}
\begin{questions}
\item $\rho=32,4/12,0=\qty{2,70}{g.cm^{-3}}$ : l’aluminium convient.
\item La mesure et les références ont des incertitudes ; un alliage peut donner une valeur voisine. C’est une identification à confirmer.
\end{questions}
\end{corrige}\end{blocLivre}
\begin{blocLivre}{cor-2-8}\begin{corrige}{8}{Du volume à la masse}
\begin{questions}
\item $m=\rho V=1,20\times40,0=\qty{48,0}{g}$.
\item Non : 21 \% est une proportion volumique. Les gaz n’ont pas tous la même masse volumique.
\end{questions}
\end{corrige}\end{blocLivre}
\begin{blocLivre}{cor-2-9}\begin{corrige}{9}{Un bilan de mélange}
\begin{questions}
\item $m_s=0,10\times40=\qty{4}{g}$.
\item La masse totale vaut 100 g ; le pourcentage final est $100\times4/100=4\,\%$.
\end{questions}
\end{corrige}\end{blocLivre}
\clearpage
\begin{blocLivre}{cor-2-10}\begin{corrige}{10}{Une pureté à discuter}
\begin{questions}
\item Ces indices sont compatibles avec l’éthanol, mais ne prouvent pas la pureté ni l’absence d’autres espèces.
\item Étudier le palier de température à pression contrôlée, compléter par une chromatographie adaptée ou un autre test pertinent.
\end{questions}
\end{corrige}\end{blocLivre}
\begin{blocLivre}{cor-2-11}\begin{corrige}{11}{Préparer une composition}
\begin{questions}
\item $m_{\mathrm{Sn}}=0,120\times2,50=\qty{0,300}{kg}$ et $m_{\mathrm{autre}}=\qty{2,20}{kg}$.
\item Les fractions massiques portent sur les masses ; les fractions volumiques portent sur les volumes, reliés aux masses par des masses volumiques différentes.
\end{questions}
\end{corrige}\end{blocLivre}
\begin{blocLivre}{cor-2-12}\begin{corrige}{12}{Un objet creux}
\demarche{Séparer le volume de métal du volume extérieur de l’objet.}
\begin{questions}
\item $V_m=m/\rho=54/2,7=\qty{20}{cm^3}$.
\item $V_{\mathrm{cavité}}=30-20=\qty{10}{cm^3}$. Le volume extérieur n’est pas le volume occupé par le métal seul.
\end{questions}
\retiens{La cavité augmente le volume extérieur sans ajouter une masse de métal correspondante.}
\end{corrige}\end{blocLivre}

